\section{EPC analysis}

% ── slides p.2–p.7 ──────────────────────────────────────────
EPCs are graphically clean but semantically loose: the XOR-join is disputed when several inputs arrive (fire once per input? block?), and for the OR-join \emph{no} semantics in the literature works on every diagram: every published definition fails somewhere.\footnote{\emph{Petri Net Transformations for Business Processes: A Survey}, §4.2; Weske M., \emph{Business Process Management: Concepts, Languages, Architectures}, Springer, Ch.~6.} Multiple start and end events are allowed, but the diagram does not record which combinations are intended. The strategy is to \emph{translate} the EPC into a workflow net, declare the EPC \textbf{sound} iff the net is, and reuse the verification machinery. The translation is not canonical (especially around the OR-join) and the soundness verdict depends on the choices made.

% ── slides p.8–p.16 ──────────────────────────────────────────
\subsection{The translation pipeline}

Three steps, supported by WoPeD: \textbf{Step 1} replaces each EPC node by a Petri-net fragment (event $\leadsto$ place, function $\leadsto$ transition, connector $\leadsto$ a small fragment depending on kind and split/join role); \textbf{Step 2} wires the fragments along the EPC arcs, either in \emph{dummy style} (insert a dummy place/transition at each seam: verbose but mechanical) or in \emph{fusion style} (merge compatible endpoints: tighter, falls back to dummies at irregular junctions); \textbf{Step 3} enforces the workflow shape, funnelling multiple start events from a fresh \textsf{start} place (via an XOR-, AND-, or mixed split, depending on whether the starts are exclusive, joint, or independent: information the EPC does not record) and dually funnelling the sinks into a fresh \textsf{end} place. Three strategies instantiate the pipeline:
\begin{center}
\small
\begin{tabular}{clll}
\toprule
 & \textbf{ingenuity / style} & \textbf{applicability} & \textbf{outcome} \\
\midrule
1st & easy, fusion & any EPC & likely unsound; relaxed soundness \\
2nd & medium, direct (dummy opt.) & no OR, strict alternation & free-choice net, fast analysis \\
3rd & hard, dummy & decorated EPC (pairings, policies) & accurate analysis \\
\bottomrule
\end{tabular}
\end{center}

% ── slides p.17–p.25 ──────────────────────────────────────────
\subsection{First approach: straight translation and relaxed soundness}

Following Dehnert and Rittgen,\footnote{Dehnert J., Rittgen P., ``Relaxed Soundness of Business Processes'', \emph{CAiSE 2001}, LNCS 2068, Springer.} each connector is replaced by a fragment that preserves \emph{all} possibly intended behaviours: EPCs are agreed upon without committing to a semantics, and the translation keeps every interpretation alive, judging the result under a weaker criterion. The fragments:
\begin{center}
\begin{tikzpicture}[node distance=0.4cm and 0.55cm]
  \node[bpmn event] (a) {};
  \node[bpmn task, rounded corners=0pt, minimum width=0.5cm, below=of a] (at) {};
  \node[bpmn event, below left=0.3cm and 0.3cm of at] (a1) {};
  \node[bpmn event, below right=0.3cm and 0.3cm of at] (a2) {};
  \draw[bpmn flow] (a) -- (at); \draw[bpmn flow] (at) -- (a1); \draw[bpmn flow] (at) -- (a2);
  \node[font=\scriptsize, below=1.6cm of a] {AND-split};
  \begin{scope}[xshift=2.6cm]
  \node[bpmn event] (b1) {};
  \node[bpmn event, right=0.5cm of b1] (b2) {};
  \node[bpmn task, rounded corners=0pt, minimum width=0.5cm] (bt) at ($(b1)!0.5!(b2)+(0,-0.8)$) {};
  \node[bpmn event, below=0.3cm of bt] (bo) {};
  \draw[bpmn flow] (b1) -- (bt); \draw[bpmn flow] (b2) -- (bt); \draw[bpmn flow] (bt) -- (bo);
  \node[font=\scriptsize, below=1.6cm of b1, xshift=0.3cm] {AND-join};
  \end{scope}
  \begin{scope}[xshift=5.4cm]
  \node[bpmn event] (c) {};
  \node[bpmn task, rounded corners=0pt, minimum width=0.45cm, below left=0.3cm and 0.25cm of c] (c1) {};
  \node[bpmn task, rounded corners=0pt, minimum width=0.45cm, below right=0.3cm and 0.25cm of c] (c2) {};
  \node[bpmn event, below=0.3cm of c1] (co1) {};
  \node[bpmn event, below=0.3cm of c2] (co2) {};
  \draw[bpmn flow] (c) -- (c1); \draw[bpmn flow] (c) -- (c2);
  \draw[bpmn flow] (c1) -- (co1); \draw[bpmn flow] (c2) -- (co2);
  \node[font=\scriptsize, below=1.6cm of c] {XOR-split};
  \end{scope}
  \begin{scope}[xshift=8.2cm]
  \node[bpmn event] (d1) {};
  \node[bpmn event, right=0.7cm of d1] (d2) {};
  \node[bpmn task, rounded corners=0pt, minimum width=0.45cm, below=0.3cm of d1] (dt1) {};
  \node[bpmn task, rounded corners=0pt, minimum width=0.45cm, below=0.3cm of d2] (dt2) {};
  \node[bpmn event] (do) at ($(dt1)!0.5!(dt2)+(0,-0.8)$) {};
  \draw[bpmn flow] (d1) -- (dt1); \draw[bpmn flow] (d2) -- (dt2);
  \draw[bpmn flow] (dt1) -- (do); \draw[bpmn flow] (dt2) -- (do);
  \node[font=\scriptsize, below=1.6cm of d1, xshift=0.35cm] {XOR-join};
  \end{scope}
  \begin{scope}[xshift=11.4cm]
  \node[bpmn event] (e) {};
  \node[bpmn task, rounded corners=0pt, minimum width=0.4cm] (e1) at ($(e)+(-0.8,-0.8)$) {\tiny l};
  \node[bpmn task, rounded corners=0pt, minimum width=0.4cm] (e2) at ($(e)+(0,-0.8)$) {\tiny lr};
  \node[bpmn task, rounded corners=0pt, minimum width=0.4cm] (e3) at ($(e)+(0.8,-0.8)$) {\tiny r};
  \node[bpmn event] (eo1) at ($(e)+(-0.5,-1.7)$) {};
  \node[bpmn event] (eo2) at ($(e)+(0.5,-1.7)$) {};
  \draw[bpmn flow] (e) -- (e1); \draw[bpmn flow] (e) -- (e2); \draw[bpmn flow] (e) -- (e3);
  \draw[bpmn flow] (e1) -- (eo1);
  \draw[bpmn flow] (e2) -- (eo1); \draw[bpmn flow] (e2) -- (eo2);
  \draw[bpmn flow] (e3) -- (eo2);
  \node[font=\scriptsize, below=1.85cm of e] {OR-split};
  \end{scope}
\end{tikzpicture}
\end{center}
The AND fragments are deterministic (one synchronising transition). The XOR-split is the free-choice conflict on a shared place; the XOR-join lets the first arrival through. The OR connectors are the awkward ones: one transition \emph{per non-empty subset} of branches (left, right, both: growing exponentially with the arity), so that every interpretation remains executable; the OR-join mirrors the split with a subset-consuming fan-in. Fusion then merges compatible endpoints across each EPC arc (element fusion when a place meets a transition; arc fusion identifying duplicated seam places).

% ── slides p.26–p.56 ──────────────────────────────────────────
\begin{examplebox}{Worked example: goods reception}
The running EPC: \textsf{goods arrived} triggers an AND-split into a left branch: \textsf{check goods}, then an XOR (\textsf{ok} directly, or \textsf{not ok} $\to$ \textsf{complaint} $\to$ \textsf{data revised} rejoining via an XOR-join), then \textsf{store goods} producing \textsf{stored}, and a right branch through an OR-join into \textsf{record receipts of goods}, producing \textsf{goods recorded}. The translation expands every connector (the OR-join becoming its multi-transition subset fan-in), fuses the seam places, and closes Step 3 with an AND-join into a fresh \textsf{end} (the upstream AND-split guarantees both sinks fire).
\par\smallskip
\emph{Straight translation verdict}: \textbf{not sound}. WoPeD flags four unbounded places, all around the OR-join fragment (three inside it, plus \textsf{goods recorded} immediately downstream): with tokens on both incoming branches, the intended both-inputs variant is enabled; but so are the single-input variants, and firing one of those consumes one token and strands the other, which can then multiply. Proper completion fails; the net is unbounded.
\par\smallskip
\emph{Repair 1: replace the OR-join by an AND-join} (keep only the both-inputs transition): now \textbf{not sound the other way}: twelve non-live transitions. When the XOR takes the \textsf{ok} arm, the branch that should feed the join's second input never fires, and the AND-join deadlocks. Unboundedness has been traded for deadlock.
\par\smallskip
\emph{Repair 2: AND-join plus an ad-hoc arc} feeding the missing input from the \textsf{ok} arm: WoPeD reports \textbf{sound}, all panels green. But the patch lives in the net world: the result is no longer the translation of any EPC, the audit trail to the modeler is broken, and reverse-engineering an EPC that translates to the patched net yields a diagram more complex and less readable than the original: whose analysis would have to restart from scratch. The first approach is a quick diagnostic, not a repair loop.
\end{examplebox}

% ── slides p.57–p.67 ──────────────────────────────────────────
\subsection{Relaxed soundness}

Strict soundness is too demanding for early-design EPCs, which legitimately contain under-specified control flow. The \emph{sound behaviours} of a net are its complete runs, $L(N) = \{\sigma \mid i \xrightarrow{\sigma} o\}$; unsound runs (deadlocks, leaked tokens) are not garbage: they pinpoint candidate mistakes.

\begin{definitionbox}{Relaxed soundness}
A workflow net $N$ is \textbf{relaxed sound} iff every transition participates in at least one sound run:
\[
  \forall t \in T.\ \exists \sigma \in L(N).\ \vec{\sigma}(t) > 0.
\]
Bad runs may exist, but no transition is structurally dead with respect to sound behaviour.
\end{definitionbox}

The goods-reception net (straight translation) is relaxed sound: one sound run takes the \textsf{not ok} arm and the both-inputs OR variant, another takes the \textsf{ok} arm and a single-input variant, and together they cover every transition. The AND-join repair is \emph{not} relaxed sound as a net (one fragment transition joins no sound run) yet it \emph{is} relaxed sound \emph{as an EPC}: every event, function, and connector appears in some sound run, since several net transitions collapse into one EPC connector. Granularity matters. Two caveats close the topic: relaxed soundness is verified by \emph{enumerating} sound runs until all transitions are covered, and (open problem) no cheaper structural characterisation is known.

% ── slides p.68–p.83 ──────────────────────────────────────────
\subsection{Second approach: simplified EPCs and free-choice nets}

Following van der Aalst,\footnote{van der Aalst W.M.P., ``Formalization and Verification of Event-driven Process Chains'', Eindhoven University of Technology, 1999.} restrict the language so the translation lands in the \emph{free-choice} class, where soundness is decidable in polynomial time. Two syntactic constraints: strict \textbf{event/function alternation} along every path (enforced beforehand by a Step~0 that inserts dummy events and functions on broken arcs, including on direct connector-to-connector arcs), and \textbf{no OR connectors} at all.

The connector translation becomes \emph{context-dependent}: the surrounding nodes already provide places (events) or transitions (functions), and only the missing kind is materialised.
\begin{itemize}
  \item \textbf{AND-split / AND-join}: between a function and events (or events and a function), the function's transition simply gains multiple output (input) arcs: no new nodes; between an event and functions, a fresh synchronising transition is inserted after the place (dually for the join).
  \item \textbf{XOR-split}: out of an event, the place feeds the competing function-transitions directly: the free-choice conflict; out of a function, a fresh place plus one fresh transition per branch are inserted.
  \item \textbf{XOR-join}: into an event, the branch transitions feed the place directly; into a function, fresh per-branch transitions funnel into a fresh place feeding it.
\end{itemize}
The uniform principle: two adjacent events need an intermediate transition, two adjacent functions an intermediate place. The result is always a free-choice workflow net. Free choice buys fast analysis, not correctness: the farmhouse-visit EPC, normalised into the sub-class, still translates to a net that WoPeD rejects: two places of the AND-join fragment (its internal place and a dummy) left uncovered by any S-component, ten non-live transitions, a reachable deadlock. The structural defect of the diagram survives the translation.

% ── slides p.84–p.119 ──────────────────────────────────────────
\subsection{Third approach: decorated EPCs}

Following Rittgen,\footnote{Rittgen P., ``Modified EPCs and Their Formal Semantics'', Universit\"at Koblenz-Landau, 1999.} keep the full connector alphabet (OR included) but require the designer to \emph{decorate} the diagram, as introduced in the EPC chapter: every (X)OR-join is paired with a \textbf{corresponding split}, and every OR-join carries a \textbf{policy}: \textsf{wfa} (wait-for-all: synchronise every activated branch; the safe default, mandatory when the corresponding split matches in type), \textsf{et} (every-time: each arriving token triggers the output independently), or \textsf{fc} (first-come: first arrival triggers, later ones are ignored). XOR-joins have no policy: their semantics (block if both inputs activate, fire on the unique one) is guaranteed by the required matching split (the implicit start split counts).

Step 1 then picks the join fragment by \emph{policy and corresponding-split type}: with a matching OR-split, the subset choice is routed from split to join, which consumes exactly the activated subset; \textsf{wfa} under an AND-split degenerates to a plain AND-join; \textsf{wfa}/\textsf{fc}/\textsf{et} under an XOR-split all coincide (only one input ever arrives), realised by two routing lanes carrying the split's commitment to the join. The dangerous combinations are explicit in the fragments: \textsf{et} under an AND-split lets \emph{both} tokens through: multiple tokens in the target, safeness broken by design; \textsf{fc} under an AND-split lets one through and \emph{strands} the other: pending tokens by design. Step 2 proceeds dummy-style (direct conversion on compatible seams; dummy transition between two places, dummy place between two transitions).

On the farmhouse example (upper OR-join decorated \textsf{fc} with the XOR-split as corresponding split, lower OR-join \textsf{wfa} with the AND-split) the fully determined translation still yields a net that WoPeD rejects: a free-choice violation, two places outside every S-component, thirteen non-live transitions. The verdict is the deepest of the three: the diagram itself contains a genuine semantic mismatch between its XOR choice and its AND synchronisation, and \emph{no} policy can paper over it. The modeler must revisit the EPC.

% ── slides p.120–p.123 ──────────────────────────────────────────
\begin{notebox}{EPC analysis: pros and cons of the three approaches}
\emph{Leave complete freedom} (1st): accepts any EPC, but most are unsound and the diagnosis is shallow: relaxed soundness only. \emph{Constrain the diagrams} (2nd): polynomial-time precision on the free-choice image, but people like flexible syntax and ignore guidelines. \emph{Require decorations} (3rd): the most accurate analysis, but people are lazy: decorations get skipped or misread. There is no universally good answer; each approach is a point on the trade-off between expressiveness, precision, and modeler discipline.
\end{notebox}

\begin{examplebox}{Exercise: customer-order EPC}
A customer-order process (compose order data, accept/reject, check availability, produce if needed, ship, check payment, complete) uses only XOR and AND connectors. It fits the simplified sub-class, so the second technique applies: fusion-style translation, Step 3 closing the two mutually exclusive end events (\textsf{rejected}, \textsf{completed}) with an XOR end pattern. The resulting free-choice workflow net passes every WoPeD check, two S-components covering all places, well-structured handles, bounded, live: \textbf{sound}. When the EPC fits the sub-class, the second technique is decisive and fast.
\end{examplebox}
